% Abstract and Section I as in the shared Overleaf main.tex, 2026-09-14 (mentor-owned).
\begin{abstract}

Vision-Language-Action (VLA) models have shown strong performance in robotic manipulation, but their high inference cost can limit practical deployment.
Recent training-free methods aim to improve efficiency without modifying or retraining the underlying policy.
However, existing studies evaluate these methods on only a small number of models or benchmarks, so their behavior across different VLA systems has not been studied in a consistent way.
In this work, we present a systematic study of training-free VLA efficiency methods across three parts of the control loop: \emph{what} the policy observes, \emph{when} it is evaluated, and \emph{how much} computation is used at each evaluation.
We study five representative methods (tricks) -- visual foveation, action repetition, depth pruning, guarded action reuse, and temporal feature fusion.
We evaluate these tricks under a common protocol across seven VLA backbones and multiple manipulation environments, measuring both task performance and efficiency.
Our results show that no single method consistently improves both efficiency and task performance across all models and settings. Different methods work better for different VLA backbones and tasks, and the same method can produce different outcomes across evaluation settings.
These findings show the importance of evaluating efficiency methods across a broader range of VLA systems rather than relying on results from a single model or benchmark.
We hope this study provides a useful empirical basis for the design and evaluation of future efficient VLA systems.

\end{abstract}


\section{INTRODUCTION}

Vision-language-action (VLA) models are emerging as general-purpose policies for instruction-conditioned robotic manipulation~\cite{openvla,spatialvla,univla}. However, their growing capabilities come with substantial inference cost. Large vision encoders and language backbones repeatedly process high-dimensional observations, while some policies additionally rely on iterative action-generation modules. In closed-loop manipulation, this cost directly affects how frequently the policy can observe the environment, update its decisions, and respond to changes during execution. Improving VLA efficiency without modifying or retraining the underlying policy is therefore important for practical deployment.

Recent training-free approaches reduce inference cost by exploiting redundancy at different stages of the VLA pipeline. Visual-token methods remove or reuse redundant image features~\cite{fastv,vlacache,vlapruner}; depth-pruning methods skip transformer layers whose representations change little~\cite{shortgpt,efficientvla}; and temporal methods reuse computation or actions across nearby control steps~\cite{vlacache,flashvla,ttfvla}. These methods often demonstrate favorable trade-offs between reduced computation and preserved task performance. However, such gains are typically reported for a specific intervention, backbone, benchmark, and implementation. This leaves an important question for efficient VLA deployment: \textbf{\emph{do training-free efficiency interventions provide consistent benefits across different VLA models and manipulation settings, or are their gains inherently configuration-dependent?}}

This question is particularly important for VLAs because efficiency interventions operate within a closed-loop control system. Unlike compression of a static recognition model, modifying VLA inference can change both the computational cost and the behavior of the policy. For example, action repetition reduces the number of policy evaluations but also reduces the frequency of environmental feedback. Visual foveation changes the information presented to the policy without necessarily reducing the number of visual tokens. Temporal reuse introduces information from previous observations, while depth pruning may affect different functions across different VLA architectures. Consequently, reduced computation does not necessarily lead to lower runtime or preserved manipulation performance. The effectiveness of an intervention (tricks) can therefore depend on the underlying architecture, action representation, task, embodiment, and implementation.

In this work, we investigate this question through a study of training-free VLA efficiency interventions. We examine when and where commonly used interventions remain effective across models and manipulation environments. We organize the interventions along three complementary dimensions of the VLA control loop: \textbf{\emph{what}} visual information the policy observes, \textbf{\emph{when}} the policy is evaluated, and \textbf{\emph{how much}} computation is performed during each evaluation. Specifically, we evaluate visual foveation, action repetition, depth pruning, guarded action reuse, and temporal feature fusion across \textit{\textbf{seven}} VLA backbones and \textbf{\textit{three}} sets for benchmarks. All interventions are evaluated against the corresponding unmodified policy using matched initial conditions, identical task instances, and a common latency protocol. We additionally retain per-episode outcomes to analyze how individual episodes change under each intervention rather than relying only on aggregate success rates.

Our results show that the effectiveness of training-free VLA efficiency interventions is highly configuration-dependent. No single method consistently improves both efficiency and manipulation performance across all models and settings, with outcomes varying across backbones and implementation choices.

Our main contributions and observations include the following: 1) \textbf{Evaluation of Key Efficiency Tricks:} We systematically evaluate five representative training-free efficiency tricks -- visual foveation, action repetition, depth pruning, guarded action reuse, and temporal feature fusion -- that target different sources of inference cost in VLA systems. Key insights include: no single trick consistently improves both efficiency and manipulation performance across all settings; the effectiveness of visual and temporal interventions varies substantially across VLA backbones and tasks; and reducing computation does not necessarily translate into proportional runtime gains or preserved task success. 2) \textbf{Cross-Model and Implementation Analysis:} We further study how the behavior of these tricks changes across model architectures and implementation choices. Our empirical results show that the same intervention can improve one backbone while degrading another under comparable evaluation conditions. 3) \textbf{Comprehensive Benchmarking:} We conduct extensive experiments across seven representative VLA backbones and multiple manipulation environments under a unified evaluation protocol, measuring both task performance and inference efficiency. Our results provide a common empirical benchmark for understanding when training-free efficiency interventions transfer across VLA systems.
